\documentclass[11pt,a4paper]{article}

\usepackage[a4paper,top=18mm,bottom=19mm,right=19mm,left=19mm]{geometry}
\usepackage{amsmath,amssymb}
\usepackage{xcolor}
\usepackage{booktabs}
\usepackage{array}
\usepackage{enumitem}
\usepackage[explicit]{titlesec}
\usepackage[most]{tcolorbox}
\usepackage{fancyhdr}
\usepackage[colorlinks=true,linkcolor=blue!45!black]{hyperref}
\usepackage{xepersian}

\settextfont[
  Path={../booklet/font/},
  UprightFont={IRLotusICEE.ttf},
  BoldFont={IRLotusICEE_Bold.ttf},
  BoldItalicFont={IRLotusICEE_BoldIranic.ttf},
  ItalicFont={IRLotusICEE_Iranic.ttf},
  FontFace={b}{n}{IRLotusICEE_Bold.ttf},
  FontFace={bx}{n}{IRLotusICEE_Bold.ttf},
  FontFace={b}{it}{IRLotusICEE_BoldIranic.ttf},
  FontFace={bx}{it}{IRLotusICEE_BoldIranic.ttf},
  FontFace={m}{sl}{IRLotusICEE_Iranic.ttf},
  Scale=1.15
]{IRLotusICEE.ttf}
\setlatintextfont[
  Path={../booklet/font/},
  BoldFont={LiberationSerif-Bold.ttf},
  BoldItalicFont={LiberationSerif-BoldItalic.ttf},
  ItalicFont={LiberationSerif-Italic.ttf}
]{LiberationSerif-Regular.ttf}
\setdigitfont[Path={../booklet/font/},Scale=1.15]{IRLotusICEE.ttf}
\defpersianfont\titlefont[Path={../booklet/font/},Scale=1]{IRTitr.ttf}

\definecolor{navy}{RGB}{20,65,105}
\definecolor{blueback}{RGB}{237,245,252}
\definecolor{greenback}{RGB}{238,248,241}
\definecolor{yellowback}{RGB}{255,248,225}
\definecolor{redback}{RGB}{254,239,239}

\newcommand{\eng}[1]{\lr{#1}}
\newcommand{\term}[2]{\textbf{#1} \eng{(#2)}}
\newcommand{\key}[1]{\textcolor{navy}{\textbf{#1}}}
\newcommand{\RPE}{\mathrm{RPE}}

\newtcolorbox{recall}[1]{
  enhanced,breakable,colback=blueback,colframe=navy,title={#1},
  fonttitle=\bfseries,boxrule=.6pt,arc=1.5mm,
  right=3.5mm,left=3.5mm,top=2mm,bottom=2mm,
  before skip=7pt,after skip=7pt,
  boxed title style={left=2mm,right=2mm,top=.7mm,bottom=.7mm}
}
\newtcolorbox{formula}[1]{
  enhanced,breakable,colback=greenback,colframe=green!45!black,title={#1},
  fonttitle=\bfseries,boxrule=.6pt,arc=1.5mm,
  right=3.5mm,left=3.5mm,top=2mm,bottom=2mm,
  before skip=7pt,after skip=7pt,
  boxed title style={left=2mm,right=2mm,top=.7mm,bottom=.7mm}
}
\newtcolorbox{trap}[1]{
  enhanced,breakable,colback=yellowback,colframe=orange!65!black,title={#1},
  fonttitle=\bfseries,boxrule=.6pt,arc=1.5mm,
  right=3.5mm,left=3.5mm,top=2mm,bottom=2mm,
  before skip=7pt,after skip=7pt
}

\setlist[itemize,1]{
  label=\textcolor{navy}{\large\textbullet},
  rightmargin=7mm,leftmargin=2mm,topsep=4pt,itemsep=3.5pt,parsep=1pt
}
\setlist[itemize,2]{
  label=\textcolor{navy}{--},
  rightmargin=7mm,leftmargin=2mm,topsep=2pt,itemsep=2.5pt,parsep=0pt
}
\setlist[enumerate]{
  label=\textcolor{navy}{\arabic*.},
  rightmargin=9mm,leftmargin=2mm,topsep=4pt,itemsep=3pt
}
\setlength{\parindent}{0pt}
\setlength{\parskip}{3pt}
\setlength{\fboxsep}{5pt}
\linespread{1.12}
\renewcommand{\arraystretch}{1.25}

\titleformat{\section}[block]
  {\normalfont\titlefont\Large\color{white}}
  {}
  {0pt}
  {\colorbox{navy}{\parbox{\dimexpr\textwidth-2\fboxsep\relax}{\raggedleft\strut#1\strut}}}
\titlespacing*{\section}{0pt}{13pt}{10pt}

\titleformat{\subsection}[hang]
  {\normalfont\bfseries\large\color{navy}}
  {\thesubsection}
  {7pt}
  {#1}
  [\vspace{2pt}\color{navy!35}\titlerule]
\titlespacing*{\subsection}{0pt}{11pt}{7pt}

\pagestyle{fancy}
\fancyhf{}
\fancyhead[R]{\small مرور فصل ۵: پاداش، یادگیری تقویتی و تصمیم‌گیری}
\fancyhead[L]{\small علوم شناختی محاسباتی}
\fancyfoot[C]{\thepage}
\setlength{\headheight}{14pt}

\begin{document}

\begin{center}
  {\titlefont\LARGE خلاصهٔ امتحانی فصل ۵}\\[3pt]
  {\Large پاداش، یادگیری تقویتی و تصمیم‌گیری}\\[3pt]
  {\small مرور تفکیک‌شدهٔ مفاهیم، فهرست‌ها و فرمول‌ها}
\end{center}

\section{بخش اول: سیستم پاداش و یادگیری تقویتی}

\subsection{مفاهیم پایهٔ تعامل با محیط}

\begin{itemize}
  \item \term{عامل}{Agent}، \term{وضعیت}{State}، \term{عمل}{Action} و \term{گذار}{Transition}:
  عامل در یک وضعیت قرار دارد، عملی انجام می‌دهد و از طریق گذار به وضعیت بعدی می‌رود:
  \[
  S\xrightarrow{A}S'
  \]

  \item \term{تقویت‌کننده}{Reinforcer}: پاداش یا تنبیهی که باعث یادگیری، تکرار یا کاهش یک رفتار می‌شود.

  \item \term{ماژول ارزش‌گذاری}{Utility Module}: کیفیت عمل انجام‌شده را بررسی می‌کند تا مدل قبلی یا \eng{Prior Model} به‌روزرسانی شود.
\end{itemize}

\subsection{مراحل معماری شناختی در مواجهه با محرک}

\begin{enumerate}
  \item دریافت محرک‌های حسی از طریق ماژول‌های بینایی یا شنیداری.
  \item انتخاب بخشی از اطلاعات محیط با فرایند توجه انتخابی.
  \item ورود اطلاعات فعلی به حافظهٔ کوتاه‌مدت.
  \item بازیابی اطلاعات مرتبط از حافظهٔ بلندمدت.
  \item ترکیب اطلاعات فعلی و دانش قبلی برای استخراج الگو.
  \item انتخاب عمل و اجرای آن از مسیر ماژول حرکتی.
  \item دریافت پاداش یا تنبیه در وضعیت بعدی و اصلاح مدل قبلی.
\end{enumerate}

\subsection{پاداش در مغز و دوپامین}

\begin{itemize}
  \item \term{دوپامین}{Dopamine}: در واکنش به یک تغییر حالت معنادار و مهم ترشح می‌شود؛ بنابراین صرفاً نشانهٔ «لذت» نیست.

  \item هر تجربهٔ خوب الزاماً دوپامین زیادی ایجاد نمی‌کند؛ اگر تجربه کاملاً قابل پیش‌بینی باشد، غافلگیری و پاسخ دوپامینی کمتر است.

  \item \term{ناحیهٔ تگمنتوم شکمی}{Ventral Tegmental Area; VTA}: ناحیهٔ مهم تولید دوپامین و محاسبهٔ خطای پیش‌بینی پاداش.
\end{itemize}

\subsection{مسیرهای دوپامینرژیک مرتبط با \eng{VTA}}

\[
\boxed{
\text{\lr{VTA}}\rightarrow
\begin{cases}
\text{\lr{Mesolimbic Pathway}}\\
\text{\lr{Mesocortical Pathway}}
\end{cases}}
\]

\begin{itemize}
  \item \term{مسیر مزولیمبیک}{Mesolimbic Pathway}:
  \begin{itemize}
    \item مسیر: \eng{VTA} به سیستم لیمبیک، به‌ویژه \eng{Nucleus Accumbens}.
    \item مسئول احساس رضایت، شرطی‌شدن، تقویت رفتار، اعتیاد و میل به تکرار رفتار.
    \item ارزش محرک یا عمل را در حافظهٔ بلندمدت تقویت یا تضعیف می‌کند.
    \item استعاره: \key{گاز}؛ «این رفتار قبلاً خوب بوده است، دوباره انجامش بده.»
  \end{itemize}

  \item \term{مسیر مزوکورتیکال}{Mesocortical Pathway}:
  \begin{itemize}
    \item مسیر: \eng{VTA} به قشر پیش‌پیشانی یا \eng{Prefrontal Cortex}.
    \item مسئول تصمیم‌گیری، تحلیل، کنترل رفتار و ارزیابی پیامدهای بلندمدت.
    \item پاداش فوری را در برابر هزینه و نتیجهٔ آینده بررسی می‌کند.
    \item استعاره: \key{ترمز یا تحلیل‌گر}.
  \end{itemize}
\end{itemize}

\subsection{اختلال‌ها و پدیده‌های مرتبط با سیستم پاداش}

\begin{itemize}
  \item \textbf{پیری:}
  \begin{itemize}
    \item تجربهٔ بیشتر \(\Leftarrow\) پیش‌بینی‌پذیری بیشتر محرک‌ها.
    \item نزدیک‌شدن خطای پیش‌بینی به صفر و کاهش ترشح دوپامین.
    \item کاهش به‌روزرسانی مدل و میل به تجربهٔ چیزهای تازه یا \eng{Exploration}.
    \item تجربه‌های تازه و کمتر تجربه‌شده می‌توانند غافلگیری سالم ایجاد کنند.
  \end{itemize}

  \item \textbf{افسردگی:}
  \begin{itemize}
    \item سیستم پاداش به‌درستی فعال نمی‌شود و کارهای قبلاً لذت‌بخش جذابیت خود را از دست می‌دهند.
    \item فرد ممکن است با عمل‌نکردن و تجربه‌نکردن در حالت قطعیت باقی بماند.
    \item خطای نزدیک صفر در این حالت ناسالم است؛ زیرا از تجربه و دقیق‌ترشدن مدل حاصل نشده است.
  \end{itemize}

  \item \textbf{اعتیاد:}
  \begin{itemize}
    \item تجربهٔ اولیه می‌تواند خطای پیش‌بینی مثبت ایجاد کند:
    \[
    \RPE>0
    \]
    \item با تکرار و سازگاری مغز، تجربه عادی می‌شود:
    \[
    \RPE\rightarrow0
    \]
    \item فرد برای ایجاد غافلگیری مثبت دوباره، دوز را افزایش می‌دهد.
    \item مغز در واکنش دفاعی، حساسیت به دوپامین را کاهش می‌دهد؛ در نتیجه لذت از محرک اعتیادآور و محرک‌های عادی کمتر می‌شود.
    \item مقاومت در برابر این چرخه به میزان مواجهه، آگاهی، محیط، ژنتیک و قدرت مسیر مزوکورتیکال بستگی دارد.
    \item \eng{Marshmallow Test}: نمونه‌ای از خودکنترلی، تأخیر پاداش و توجه به نتیجهٔ بلندمدت.
  \end{itemize}
\end{itemize}

\begin{formula}{چرخهٔ اعتیاد}
\[
\text{تجربهٔ قوی}\rightarrow\RPE>0
\rightarrow\text{سازگاری}\rightarrow\RPE\approx0
\rightarrow\text{افزایش دوز}\rightarrow\RPE>0\rightarrow\cdots
\]
\end{formula}

\subsection{کدگذاری پیش‌بینانه}

\begin{itemize}
  \item \term{کدگذاری پیش‌بینانه}{Predictive Coding}: مقایسهٔ مداوم میان دانش قبلی یا \eng{Prior Knowledge} و واقعیت مشاهده‌شده.
  \item مبنای یادگیری، خود واقعیت خام نیست؛ بلکه \key{خطای میان پیش‌بینی و واقعیت} است.
  \item در سیستم پاداش نیز انتظار قبلی با پاداش واقعی مقایسه می‌شود.
\end{itemize}

\subsection{خطای پیش‌بینی پاداش}

\begin{formula}{Reward Prediction Error}
\[
\text{\lr{Reward Prediction Error}}
=\text{\lr{Actual Reward}}-\text{\lr{Predicted Reward}}
\]
\[
\boxed{\RPE=R_{\mathrm{actual}}-R_{\mathrm{predicted}}}
\]
\end{formula}

\begin{itemize}
  \item \(\RPE>0\): نتیجه بهتر از انتظار؛ غافلگیری مثبت و تقویت عمل یا محرک.
  \item \(\RPE=0\): نتیجه مطابق انتظار؛ اطلاعات تازه‌ای برای تغییر مدل وجود ندارد.
  \item \(\RPE<0\): نتیجه بدتر از انتظار؛ ارزش عمل یا محرک کاهش می‌یابد.
\end{itemize}

\subsection{اصل انرژی آزاد}

\begin{itemize}
  \item \term{اصل انرژی آزاد}{Free Energy Principle}: سیستم شناختی برای کاهش عدم قطعیت، غافلگیری و خطای پیش‌بینی محیط را تجربه می‌کند.
  \item هدف، هماهنگ‌شدن مدل ذهنی با واقعیت است:
  \[
  \boxed{\RPE\rightarrow0}
  \]
  \item صفرشدن سالم خطا حاصل تجربه و اصلاح مدل است؛ نه عمل‌نکردن و تجربه‌نکردن.
\end{itemize}

\section{بخش دوم: مدل‌های محاسباتی یادگیری تقویتی}

\subsection{یادگیری تقویتی و ویژگی‌های آن}

\begin{itemize}
  \item \term{یادگیری تقویتی}{Reinforcement Learning}: یادگیری از طریق تعامل فعال با محیط، انجام عمل و دریافت پاداش یا تنبیه.
  \item یادگیری معمولاً \term{برخط}{Online} و نمونه‌به‌نمونه است؛ عامل از ابتدا مدل دقیق محیط را ندارد.
  \item \textbf{در برابر یادگیری نظارت‌شده:} برچسب درست یا \eng{Label} آماده وجود ندارد.
  \item \textbf{در برابر یادگیری بدون‌ناظر:} هدف فقط کشف ساختار داده نیست؛ عمل و بازخورد پاداش بخش اصلی مسئله‌اند.
\end{itemize}

\subsection{محیط قطعی و تصادفی}

\begin{itemize}
  \item \term{محیط قطعی}{Deterministic}: یک عمل مشخص در یک وضعیت مشخص، همیشه به نتیجهٔ یکسان می‌رسد.
  \item \term{محیط تصادفی}{Stochastic}: انجام یک عمل مشخص لزوماً به یک نتیجهٔ قطعی ختم نمی‌شود و گذار با احتمال همراه است:
  \[
  S\xrightarrow{A}S'
  \]
  \item عامل باید با تجربه، دریافت بازخورد و اصلاح تدریجی مدل گذار با این عدم قطعیت کنار بیاید.
\end{itemize}

\subsection{یادگیری تفاوت زمانی}

\begin{itemize}
  \item \term{یادگیری تفاوت زمانی}{Temporal Difference Learning}: اختلاف میان ارزش پیش‌بینی‌شده در زمان \(t\) و نمونهٔ حاصل از تجربه در زمان \(t+1\) را یاد می‌گیرد.
  \item \term{نرخ یادگیری}{Learning Rate} با \(\alpha\) نشان داده می‌شود:
  \[
  \boxed{0\leq\alpha\leq1}
  \]
  \item \(\alpha=0\): تجربهٔ جدید اثری ندارد و یادگیری رخ نمی‌دهد.
  \item \(\alpha=1\): فقط تجربهٔ جدید مهم است و دانش قبلی نادیده گرفته می‌شود.
  \item \(0<\alpha<1\): درون‌یابی میان دانش قبلی و تجربهٔ جدید.
\end{itemize}

\begin{formula}{فرمول به‌روزرسانی TD}
\[
\boxed{V_{t+1}=V_t+\alpha(\text{\lr{Sample}}-V_t)}
\]
\[
\text{\lr{Sample}}-V_t=\RPE
\]
\[
\boxed{V_{t+1}=V_t+\alpha\times\RPE}
\]
\end{formula}

\begin{itemize}
  \item هدف نهایی:
  \[
  \text{\lr{Sample}}=V_t
  \quad\Longrightarrow\quad
  \RPE=0
  \]
\end{itemize}

\subsection{ضریب تخفیف}

\begin{itemize}
  \item \term{ضریب تخفیف}{Discount Factor} با \(\gamma\) نشان داده می‌شود و میزان اهمیت پاداش‌های آینده را مدل می‌کند:
  \[
  \boxed{0\leq\gamma\leq1}
  \]
  \item \(\gamma=0\): فقط پاداش فوری؛ تصمیم‌گیری کاملاً کوتاه‌نگر.
  \item \(\gamma=1\): پاداش‌های آینده بدون تخفیف در نظر گرفته می‌شوند.
  \item \(0<\gamma<1\): آینده مهم است، اما پاداش‌های دورتر اثر کمتری دارند.
\end{itemize}

\begin{formula}{مجموع پاداش‌ها با ضریب تخفیف}
\[
\boxed{
r_t+\gamma r_{t+1}+\gamma^2r_{t+2}
+\gamma^3r_{t+3}+\cdots}
\]
\end{formula}

\begin{itemize}
  \item ارتباط زیستی مطرح‌شده: سروتونین بالاتر ممکن است با عجلهٔ کمتر و توان بیشتر برای درنظرگرفتن پاداش آینده مرتبط باشد؛ البته رفتار به عوامل متعدد وابسته است.
\end{itemize}

\subsection{تابع ارزش}

\begin{itemize}
  \item \term{تابع ارزش}{Value Function}: ارزش قرارداشتن در حالت \(s\) را با \(V(s)\) نشان می‌دهد.
  \item در هر گام، ارزش همان حالتی به‌روزرسانی می‌شود که عامل از آن حرکت کرده است.
\end{itemize}

\begin{formula}{فرمول کامل به‌روزرسانی تابع ارزش}
\[
\boxed{
V_{t+1}(s)=(1-\alpha)V_t(s)
+\alpha\left[r+\gamma V_t(s')\right]}
\]
\begin{itemize}
  \item \(V_t(s)\): ارزش قبلی حالت فعلی.
  \item \(r\): پاداش دریافت‌شده.
  \item \(V_t(s')\): ارزش حالت بعدی.
  \item \((1-\alpha)V_t(s)\): سهم دانش قبلی.
  \item \(\alpha[r+\gamma V_t(s')]\): سهم تجربهٔ جدید.
\end{itemize}
\end{formula}

\subsection{تابع \eng{Q}}

\begin{itemize}
  \item \term{تابع Q}{Q Function}: ارزش انجام عمل خاص \(a\) در حالت خاص \(s\).
  \item \textbf{تفاوت اصلی:} \(V(s)\) ارزش خود حالت است؛ \(Q(s,a)\) ارزش یک عمل در آن حالت.
  \item رابطهٔ میان \(V\) و \(Q\):
  \[
  \boxed{V(s)=\max_a Q(s,a)}
  \]
  \item \term{تابع احتمال گذار}{Transition Function} با \(T(s,a,s')\) نشان داده می‌شود: احتمال رسیدن از \(s\) به \(s'\) با انجام \(a\).
\end{itemize}

\begin{formula}{محاسبهٔ Q در محیط تصادفی}
\[
\boxed{
Q(s,a)=\sum_{s'}T(s,a,s')
\left[R(s,a,s')+\gamma V(s')\right]}
\]
\end{formula}

\subsection{الگوریتم \eng{Q-Learning}}

\begin{itemize}
  \item در \eng{Q-Learning} ارزش عمل‌ها در حالت‌ها مستقیماً به‌روزرسانی می‌شود؛ نه فقط ارزش خود حالت‌ها.
  \item عبارت \(\max_{a'}Q_t(s',a')\) بهترین عمل ممکن در حالت بعدی را نشان می‌دهد.
\end{itemize}

\begin{formula}{فرمول به‌روزرسانی Q-Learning}
\[
\boxed{
Q_{t+1}(s,a)=(1-\alpha)Q_t(s,a)
+\alpha\left[r+\gamma\max_{a'}Q_t(s',a')\right]}
\]
\end{formula}

\subsection{موازنهٔ اکتشاف و بهره‌برداری}

\begin{itemize}
  \item \term{بهره‌برداری}{Exploitation}: استفاده از بهترین تجربه‌های شناخته‌شدهٔ قبلی؛ انتخاب حریصانه یا \eng{Greedy}.
  \item \term{اکتشاف}{Exploration}: فرصت‌دادن به عمل‌های کمتر امتحان‌شده برای کشف گزینه‌های احتمالاً بهتر.
  \item تصمیم‌گیری مؤثر معمولاً به تعادل میان این دو نیاز دارد.
\end{itemize}

\subsection{\eng{Active RL} و \eng{Passive RL}}

\begin{itemize}
  \item \term{یادگیری تقویتی ایستا}{Passive RL}:
  \begin{itemize}
    \item تمرکز بیشتر بر بهره‌برداری و تجربه‌های قبلی.
    \item تمایل کمتر به امتحان گزینه‌های جدید.
    \item امتیاز عمل:
    \[
    \boxed{\mathrm{Score}(s,a)=Q(s,a)}
    \]
  \end{itemize}

  \item \term{یادگیری تقویتی پویا}{Active RL}:
  \begin{itemize}
    \item عدم قطعیت را نیز در انتخاب عمل وارد می‌کند.
    \item میان اکتشاف و بهره‌برداری تعادل ایجاد می‌کند.
    \item امتیاز عمل:
    \[
    \boxed{\mathrm{Score}(s,a)=Q(s,a)+\beta U(s,a)}
    \]
    \item \(U(s,a)\): میزان عدم قطعیت دربارهٔ عمل؛ برای عمل کمتر امتحان‌شده بیشتر است.
    \item \(\beta\): میزان اهمیت عدم قطعیت در تصمیم عامل.
    \item مشابه روش \term{کران اطمینان بالا}{Upper Confidence Bound; UCB}.
  \end{itemize}
\end{itemize}

\end{document}
